\documentclass[11pt]{report}

\usepackage[margin=1in]{geometry}
\usepackage{amsmath}
\usepackage{booktabs}
\usepackage[hidelinks]{hyperref}
\usepackage{microtype}

\title{Bulk-Silicon Wannier Construction and\\Controlled-Reduction Diagnosis Research Protocol\\Common-Space Diagnosis, Hopping Analysis, and Post-S110 Continuation}
\author{Project Koios}
\date{Calculation-owned staged protocol}

\begin{document}
\pagenumbering{roman}
\maketitle

\begin{abstract}
This protocol separates scientific target-space design, local numerical optimization,
arithmetic or toolchain instability, common-space rank diagnosis, independent
verification, symmetry-sector diagnosis, scientific disposition, complete hopping
characterization, support reduction, Slater--Koster parameterization, and hidden-parent
evaluation in the bulk-silicon Wannier route.  It prevents failed localization from
being mislabeled as either proof of physical infeasibility or evidence of
floating-point instability, and it prevents successful Wannier90 convergence from being
promoted directly to physical validation.

The post-S110 order is mandatory:
\begin{align*}
 \mathrm{CSD0}&\longrightarrow\mathrm{CSD1}
 \longrightarrow\mathrm{CSD1\mbox{-}V}
 \longrightarrow\mathrm{CSD2}\longrightarrow\mathrm{CSD3},\\
 &\longrightarrow\mathrm{HF0}\longrightarrow\cdots
 \longrightarrow\mathrm{HF4},\\
 &\longrightarrow\text{one authorized successor branch}.
\end{align*}
CSD1-V shares the authenticated Wannier90 parsers and parsed matrix records used by the
producer; its independence begins at the parsed-matrix boundary and covers the
scientific matrix construction and numerical diagnosis.  The full untruncated Wannier
operator is subsequently analyzed as a complete real-space hopping representation
without removing support.  Descriptive hopping locality precedes any new support
decision, and the hopping-first disposition precedes activation of a Slater--Koster
successor branch.  Direct Kohn--Sham parent values at withheld points remain an
evaluation oracle rather than an operational reduction gate.  The protocol preserves
the completed full-source characterization (S30), revealed-parent development evaluation
(S40), shell-family design (S90), shell-family evaluation (S100), full-support disposition
(S105), and blocked ten-orbital alignment attempt (S110).  Exact thresholds, execution
authority, and evidence dispositions remain calculation-owned.
\end{abstract}

\tableofcontents
\clearpage
\pagenumbering{arabic}

\chapter{Scope and authority}

The protocol applies after an authenticated plane-wave parent, ordered NSCF Bloch frame,
Wannier90 neighbor topology, and AMN/MMN/EIG interface have been established.  Reusable
software owns parsing, identity checks, bounded execution, matrix invariants, and lineage.
The calculation owns rank, projections, windows, convergence coordinates, observable
tolerances, candidate selection, and scientific acceptance.

A trial projection initializes a Wannier gauge.  It does not impose a Slater--Koster
Hamiltonian, hopping shell, or fitted parameter set.  The complete real-space hopping
representation is analyzed before support removal.  A support disposition is then made
without turning descriptive locality into an atomic-orbital parameterization.  Only
after those gates may a Slater--Koster study begin, and only after its state spaces have
been explicitly aligned.  Agreement with withheld direct-parent values is not a
prerequisite for using the full operator as the fixed controlled-study source; such
agreement controls only physical-fidelity claims about the source and its reductions.

\chapter{Identifier-to-description crosswalk}

Identifiers are provenance shorthand, not scientific descriptions.  After the following
crosswalk, later chapters may use an identifier only for the construction or retained
stage defined here.

\begin{table}[htbp]
\centering
\small
\begin{tabular}{@{}p{0.19\textwidth}p{0.74\textwidth}@{}}
\toprule
Identifier & Scientific construction \\
\midrule
$C_6$ & Gamma-origin $6\times6\times6$ mesh with 216 reciprocal points,
12 calculated parent bands, retained rank ten, outer window
$[-5.85,19.02]$~eV, frozen window $[-5.85,8.0]$~eV, and the declared
silicon $sp^3$ plus second radial-$s$ projection family. \\
$C_{6\text{-}T1}$ & The same 12-band construction with only the outer window
expanded to $[-6.0,21.0]$~eV. \\
$C_{6\text{-}B20}$ & The same mesh, retained rank, frozen window, and projection
family with the calculated parent-band inventory increased from 12 to 20, before the
final outer-window selection. \\
$C_{6\text{-}B20\text{-}W21}$ & The retained controlled-reduction source: the
20-band construction with outer window $[-6.0,21.0]$~eV and frozen window
$[-5.85,8.0]$~eV. \\
$C_8$ & A separately proposed $8\times8\times8$ construction-mesh study; not an
automatic successor. \\
$C_{12}$ & A separately proposed $12\times12\times12$ upper mesh diagnostic; not an
automatic successor. \\
\bottomrule
\end{tabular}
\caption{Construction identifiers and their complete scientific meanings.}
\end{table}

\begin{table}[htbp]
\centering
\small
\begin{tabular}{@{}p{0.19\textwidth}p{0.74\textwidth}@{}}
\toprule
Identifier & Retained stage or disposition \\
\midrule
S30 & Complete untruncated-source characterization and source-mesh reconstruction. \\
S40 & Revealed 408-point direct-parent development evaluation. \\
S90 & Complete physical-radius shell-family design for spatial truncation. \\
S100 & Complete-shell evaluation relative to the full untruncated source. \\
S105 & Scientific disposition retaining all support, so
$H_{\mathrm{T}}=H_{\mathrm{W}}$. \\
S110 & Blocked ten-orbital trial-to-Wannier common-space alignment; no fit executed and
no Slater--Koster parameter vector produced. \\
\bottomrule
\end{tabular}
\caption{Retained stage identifiers and their scientific meanings.}
\end{table}

\chapter{Distinct failure classes}

The following classes shall not be conflated.

\section{Target-space or construction failure}

Rank, parent-band inventory, outer and frozen windows, projections, reciprocal mesh, and
neighbor topology define the represented state space and initialization.  Sensitivity to
these coordinates is scientific-construction evidence.  A construction can create a
poorly conditioned or strongly nonconvex optimization landscape without implying
arithmetic failure.

\section{Local-optimizer failure}

Wannier90 uses a projection-initialized local optimization; it does not provide a
certificate of the global minimum of the spread functional.  Iteration exhaustion,
large late spread excursions, or materially different gauges from different optimizer
coordinates establish nonstationarity or local-optimizer sensitivity.  They do not prove
that no acceptable gauge exists.

\section{Arithmetic or toolchain instability}

Arithmetic or toolchain instability requires direct evidence such as non-finite values,
corrupt or inconsistent artifacts, failed matrix identities beyond declared tolerances,
same-input nondeterminism, or controlled sensitivity to arithmetic library or executable
choice.  Optimizer excursions alone are insufficient for this diagnosis.

\section{Hidden-parent discrepancy and physical inadequacy}

A converged and structurally valid Wannier occurrence can still disagree materially
with direct Kohn--Sham parent values.  Interpolation, valley identity and location, edge
energy, derivatives, curvature or effective mass, and real-space locality are separate
evaluation channels.  At withheld points these discrepancies are known only because the
controlled study evaluates the parent model.  They are oracle labels, not diagnostics
available to a deployed reduction, and shall not be used as convergence tests,
truncation-selection rules, fitting objectives, or automatic successor triggers.
Failure of a parent-fidelity target limits physical claims but does not invalidate a
structurally eligible operator as the fixed source in a controlled reduction study.

\chapter{Mandatory diagnostic procedure}

For every localization occurrence:

\begin{enumerate}
  \item Require exact executable and input identities, a bounded authorization record,
    a fresh or explicitly declared restart state, and no silent retry.
  \item Assess disentanglement and Wannierisation separately using their native
    convergence messages and criteria.  A zero return code is not a substitute.
  \item Verify finite dimensions and k-point ordering, gauge unitarity,
    disentanglement isometry when applicable, Hamiltonian representative pairing, and
    Hermiticity.
  \item Inspect convergence trajectories.  Distinguish smooth contraction toward a
    declared threshold from intermittent plateaus, renewed excursions, and budget
    exhaustion.
  \item Hold authenticated interface artifacts fixed when diagnosing optimizer or
    target-space changes.  Record every changed scientific and algorithmic coordinate.
  \item Permit only a prospectively bounded optimizer study.  If the trajectory remains
    irregular and observably gauge-sensitive, revise the scientific construction rather
    than performing an open-ended optimizer search.
  \item After native convergence, qualify the full operator using checks available
    without withheld parent values.  Record direct-parent comparisons separately as
    development-oracle or confirmation-oracle evidence.
\end{enumerate}

\chapter{Controlled-reduction and hidden-oracle policy}

Let $H_{\mathrm{W}}$ denote the full untruncated Wannier operator,
$H_{\mathrm{T}}$ a spatially truncated operator, $H_{\mathrm{SK}}$ a declared
Slater--Koster parameterization, and $H_{\mathrm{KS}}$ the direct Kohn--Sham parent on
an aligned comparison space.  The controlled reduction studies
\begin{equation}
  H_{\mathrm{W}}\longrightarrow H_{\mathrm{T}}
  \longrightarrow H_{\mathrm{SK}}.
\end{equation}
The study shall keep at least the following channels separate:
\begin{align}
  \text{parent-to-source discrepancy} &: H_{\mathrm{KS}}-H_{\mathrm{W}},\\
  \text{spatial-truncation error} &: H_{\mathrm{W}}-H_{\mathrm{T}},\\
  \text{parameterization error} &: H_{\mathrm{T}}-H_{\mathrm{SK}},\\
  \text{oracle end-to-end discrepancy} &: H_{\mathrm{KS}}-H_{\mathrm{SK}}.
\end{align}
These expressions are schematic until basis, gauge, state space, energy reference, and
units have been aligned.  Observable comparisons may replace operator subtraction when
no justified operator alignment exists.

Operational diagnostics are quantities available without a new direct-parent
calculation at the target point.  They include native convergence, finite dimensions,
unitarity and isometry residuals, pair Hermiticity, authenticated source reconstruction,
real-space shell and tail measures, full-to-truncated operator or spectral differences,
stability across prospectively declared truncations and stencils, and fit residuals to
the declared reduction source.  Direct-parent discrepancies are evaluation overlays.
They do not calibrate, validate, or serve as targets for the operational diagnostics,
and no claim is made that an operational diagnostic predicts parent error.

The reciprocal-space roles are fixed.  The 216-point C6 construction mesh establishes
source reconstruction.  The revealed 408-point geometry is the development evaluation
inventory: complete-shell errors are measured there relative to full
$H_{\mathrm{W}}$, while direct-parent comparisons are displayed only as a separate
evaluation overlay.  The frozen 402-point geometry is confirmation-only.  Neither its
geometry nor any future parent value may define or tune diagnostics, candidates,
frontiers, truncation selection, fitting, retry selection, thresholds, or numerical
comparison policy.  The disjoint 108/108 construction-grid split remains reserved for a
future separately authorized Slater--Koster study and has no role in hopping analysis.

\chapter{Global-minimum policy}

Structural eligibility for controlled reduction does not require an unsupported claim
that a global spread minimum was found.  A locally converged candidate may serve as the
fixed reduction source only if it is structurally valid, reproducible under the declared
protocol, and satisfies the declared operational eligibility checks.  This status does
not establish physical fidelity to hidden parent values.  Competing-minimum studies
become necessary when two
or more converged initializations or constructions produce materially different target
observables or locality conclusions.  Comparisons shall use gauge-invariant spectra and
within-candidate locality metrics unless an explicit alignment map exists.

The lowest final spread among different target spaces is not, by itself, a selection
rule because changing the frozen manifold changes the represented optimization problem.

\chapter{Prospective convergence policy}

For prospective $C_{6\text{-}B20\text{-}W21}$ localization, the calculation adopts
the convergence-derived tuple from the precise protocol in pinned
\texttt{aiida-wannier90-workflows} 3.1.0.  Its builder multiplies the per-atom tolerance
meta-parameters by the number of sites.  For the exact two-atom primitive cell this
gives
\begin{align}
  \texttt{conv\_tol} &= 2.0\times10^{-9}\ \text{\AA}^{2}, &
  \texttt{conv\_window} &= 6, \\
  \texttt{dis\_conv\_tol} &= 2.0\times10^{-9}, &
  \texttt{num\_cg\_steps} &= 200,
\end{align}
with \texttt{num\_iter}=40000 and \texttt{dis\_num\_iter}=40000.  The emitted
\texttt{dis\_conv\_tol} remains dimensionless; ``per atom'' describes the workflow's
scaling convention.  Both native convergence messages remain mandatory, and the
40000-step values are fixed eligibility ceilings rather than evidence of a global
minimum or scientific adequacy.

This adaptation does not adopt the precise protocol's factor-three band count,
automatic 0.1~\AA$^{-1}$ mesh rule, PBEsol pseudopotential family, projections, energy
windows, MPI policy, or restart handlers.  Historical occurrences retain the exact
policies under which they ran and are not reclassified retrospectively.  No automatic
extension beyond the adopted ceilings is permitted.

\chapter{Current silicon evidence}

The all-conduction-frozen rank-10 construction froze eight conduction states and no
valence states at every source k point.  Its localization trajectories remained
nonstationary through as many as 4000 iterations, including 277 post-1000 spread changes
larger than $0.05\ \text{\AA}^{2}$ and a final large excursion at iteration 3934.  A
500/1000 diagnostic found close source-mesh conduction spectra but material dense-grid,
valley, and curvature differences.  These observations support optimization-landscape
and target-space sensitivity, not arithmetic instability.

The successor retained the executable, rank, projections, mesh, and byte-identical
AMN, MMN, EIG, and NNKP interface artifacts, but changed the frozen interval to
$[-5.85,8.00]$~eV in the native QE energy reference.  This freezes all four valence
states and zero to two sampled low-energy conduction states per k point.  With
\texttt{dis\_num\_iter = 2000}, disentanglement converged at iteration 1290 and
Wannierisation converged at iteration 189.  The final total spread was
$19.532839107\ \text{\AA}^{2}$; retained gauge-unitarity and disentanglement-isometry
residuals were approximately $1.58\times10^{-10}$ and $1.32\times10^{-10}$, and the
maximum Hamiltonian Hermiticity residual was $0.0$~eV.

Because the same authenticated interface and executable admit convergence after the
target-space revision, the supported diagnosis is that the earlier obstruction arose
primarily from construction-induced optimization geometry.  This is not a proof that
the new occurrence is a global spread minimum or a scientifically accepted operator.

The subsequent no-calculator characterization verified source reconstruction and
recorded the candidate's interpolation, valley, curvature, and locality behavior without
assigning pass/fail status to candidate-only observations.  A prospectively frozen
408-point direct-parent NSCF then supplied development-oracle labels.  All six
candidate-to-parent valley displacements, $0.02120$--$0.03086$~\AA$^{-1}$, met the
$0.034465$~\AA$^{-1}$ target.  The candidate indirect gap was $0.52261$~eV versus the
direct-parent $0.61125$~eV, a 14.50\% relative discrepancy against the 5\% evaluation
target.  On the frozen candidate-centered primary stencils, all transverse-mass
comparisons and two of six longitudinal-mass comparisons exceeded the 8\% evaluation
target.  These values limit physical-fidelity claims and provide controlled labels for
diagnostic development.  They are not operational rejection criteria for a reduction
whose fixed source is the full Wannier operator.

\chapter{CSD0: source, state-space, and convention authentication}

\section{Process}

CSD0 shall bind the exact development request, the authenticated S30 full-source
operator record, the S105 full-support disposition, the blocked S110 result, and the
native AMN, gauge, and disentanglement artifacts.  It shall verify every source path and
SHA-256 digest before parsing, reject paths outside the declared repository roots, and
reject an existing output root.  The shared Wannier90 parsers owned by
\texttt{projectkoios-simulations} are the only admitted file-format boundary.  CSD0 does
not require, and this protocol forbids, a second implementation of those parsers.

The admitted state space has 216 Gamma-origin construction points, 20 parent bands, and
a retained dimension of ten.  The basis order, two descriptive diamond-center labels,
radial-channel convention, spin convention, geometry, units, energy reference, and
reciprocal-point order shall be explicit and exact.  Legacy one-letter center labels may
appear only while correlating immutable S110 bytes with the descriptive current basis.

\section{Mathematics}

Let $A(\mathbf{k})$ be the parsed AMN trial-column matrix,
$U_{\mathrm{dis}}(\mathbf{k})$ the parsed disentanglement isometry, and
$U(\mathbf{k})$ the parsed Wannier gauge.  CSD0 authenticates their dimensions and the
common ordered point set before any product is formed.  It also authenticates the exact
unitary transformation $T$ from the native eight $sp^3$ hybrids plus two radial
$s$ projections to the ordered two-center
$(s,p_x,p_y,p_z,s^*)$ trial basis.  Here $s^*$ is radial channel one of the $s$
component, not a principal quantum number.

The CSD0 gate is identity-level rather than tolerance-level.  Write
\begin{equation}
 \mathcal{S}_{\mathrm{CSD0}}=(h,d,K,\mathcal{B},\mathcal{G},\mathcal{U},E_0),
\end{equation}
where $h$ is the ordered source-hash ledger, $d$ the exact dimensions, $K$ the ordered
reciprocal-point set, $\mathcal{B}$ the basis definition, $\mathcal{G}$ the geometry and
gauge conventions, $\mathcal{U}$ the unit system, and $E_0$ the energy reference.  A
missing or unequal member of $\mathcal{S}_{\mathrm{CSD0}}$ stops the sequence.  It is
not repaired by filename inference, orbital count, matrix size, a compatibility alias,
or a guessed convention.

\section{Scientific interpretation and gate}

CSD0 establishes that subsequent matrices refer to the same declared physical and
represented state spaces.  It does not establish rank, fit feasibility, symmetry
content, or physical adequacy.  Passing CSD0 makes the no-calculator CSD1 reconstruction eligible for a separately
authorized analysis run; it does not itself grant execution authority.  It authorizes
no QE or Wannier90 execution, target-space change, evidence promotion, fit, or
scientific acceptance.

\chapter{CSD1: common-space rank reconstruction}

\section{Process}

CSD1 shall reconstruct the complete 216-point trial-to-Wannier overlap field from the
CSD0-authenticated parsed records.  The producer may use the existing direct singular
value decomposition route.  It shall record every point, not only deficient points, and
shall retain the maximum, next-to-minimum, and minimum singular values, resolved rank,
condition number when finite, normalized principal-angle diagnostic, and the full-rank
gate outcome.  It shall also retain the Gamma deficient-direction weights in the exact
ordered atomic trial basis.

The current calculation-owned producer implements this projector-rank reconstruction
subset without importing the S110 calculation implementation.  It is designed to
reconstruct the retained deficient-point inventory.  Its synthetic software checks do
not establish that the current producer has completed a silicon reconstruction, and
neither software passage nor a producer-only result is independent CSD1 verification
or retained scientific acceptance.

\section{Prospective amendment: conditional arithmetic-precision escalation}

Source authentication, physical state-space associations, independent rank
reconstruction, and justified arithmetic/native-input error bounds remain mandatory.
Higher precision is a conditional diagnostic, not a routine completion requirement.
When the arithmetic bounds are decisive and the independent routes agree, no
higher-precision run or completion artifact is required.  This amendment does not
relax a rank threshold, uncertainty bound, agreement criterion, or scientific gate.

Escalation is required before progression only if arithmetic-error bounds leave a
relevant rank or direction diagnosis ambiguous, or if a discrepancy between independent
routes remains plausibly arithmetic after checking source identities, basis/gauge
conventions, and implementation errors.  A source mismatch must be repaired at its
source; it is not a reason to manufacture agreement by increasing precision.  A
bounded escalation may target the affected points rather than repeat the complete
216-point field at arbitrarily high precision.  Its arithmetic budget and stopping
condition shall be declared before that follow-up; no blanket 128/256-bit campaign
is required.  Unresolved arithmetic or disagreement still stops the sequence.

Arithmetic precision shall be kept separate from native-file precision.  The same
CSD0-authenticated parsed inputs, transformation, energy reference, native-uncertainty
model, and frozen singular-value threshold shall be used.  Extra arithmetic digits
do not recover discarded source digits, establish a rounding model, or repair
uncertain physical source/subspace associations.  If native-input uncertainty rather
than arithmetic is limiting, more arithmetic alone cannot resolve it.  No retained
analysis, calculator execution, fitted map, or evidence promotion is authorized by
this prospective amendment.

\section{Mathematics}

Define the retained parent-space frame and atomic trial columns by
\begin{align}
 W(\mathbf{k}) &= U_{\mathrm{dis}}(\mathbf{k})U(\mathbf{k}),\\
 C(\mathbf{k}) &= A(\mathbf{k})T^{\dagger},
\end{align}
and define the raw common-space overlap
\begin{equation}
 B(\mathbf{k})=W(\mathbf{k})^{\dagger}C(\mathbf{k}).
 \label{eq:csd1-overlap}
\end{equation}
For
\begin{equation}
 B(\mathbf{k})=L(\mathbf{k})\Sigma(\mathbf{k})R(\mathbf{k})^{\dagger},
\end{equation}
the resolved rank is
\begin{equation}
 r_{\delta}(\mathbf{k})=
 \#\{j:\sigma_j(\mathbf{k})>\delta_{\mathrm{native}}\},
 \label{eq:csd1-resolved-rank}
\end{equation}
where $\delta_{\mathrm{native}}$ is the prospectively declared conservative perturbation
bound induced by the retained decimal precision of AMN, gauge, and disentanglement
records.  The raw singular values in Equation~\eqref{eq:csd1-resolved-rank} remain the
frozen gate.  Principal angles may be reported only after independently
orthonormalizing the two ten-column frames; raw AMN columns are not assumed orthonormal.

Under the retained native bound
$\delta_{\mathrm{native}}=5.890206739032098\times10^{-7}$, the reconstructed field has
resolved rank nine at Gamma and at the ordered point $(0,1/3,0)$.  Their retained
minimum singular values are respectively
$5.362000386100462\times10^{-13}$ and
$2.0186562353866025\times10^{-7}$.  These are retained S110 comparison values, not
new results of the current CSD1 producer or newly accepted evidence.

\section{Scientific interpretation and gate}

A rank-nine result means that the declared ten-dimensional atomic trial frame and the
retained ten-dimensional Wannier subspace do not admit a resolved unitary
identification under the frozen gate.  It is an orbital-subspace alignment failure.  It
is not established as an energy-degeneracy failure, a floating-point failure, or proof
that every possible silicon tight-binding model is impossible.  No pseudoinverse,
pointwise completion, deleted orbital, or reduced-rank fit may be manufactured inside
CSD1.  CSD1 passes forward only to CSD1-V; it does not authorize CSD2 or a fit by itself.

\chapter{CSD1-V: independent rank verification}

\section{Process and independence boundary}

CSD1-V shall consume the same CSD0-authenticated parsed matrix records.  Parser sharing
is required: independent verification begins after the file-format boundary.  The
verifier shall not import the producer's overlap, rank, singular-value, deficient-vector,
weight-aggregation, comparison, or serialization implementations.  It shall not use a
producer-generated matrix, singular value, rank, or deficient vector as a computational
input.  Producer and verifier outputs may be correlated only after the verifier has
completed its own result.

The verifier shall use a separately implemented Gram/eigen route and write a distinct,
fail-closed verification record under a nonexistent output root.  Repository tests shall
forbid imports from \texttt{common\_space.diagnostics} and the S110 definition-fit
action while permitting the shared provider parsers and immutable parsed record types.

\section{Mathematics}

Using an independently assembled $\widetilde B(\mathbf{k})$ from the shared parsed
records, define
\begin{equation}
 G(\mathbf{k})=\widetilde B(\mathbf{k})^{\dagger}\widetilde B(\mathbf{k}).
 \label{eq:csd1v-gram}
\end{equation}
A Hermitian eigendecomposition
\begin{equation}
 G(\mathbf{k})V(\mathbf{k})=
 V(\mathbf{k})\Lambda(\mathbf{k})
\end{equation}
provides independently reconstructed singular values
\begin{equation}
 \widetilde\sigma_j(\mathbf{k})=
 \sqrt{\max(\lambda_j(\mathbf{k}),0)}.
\end{equation}
Forming a Gram matrix squares the singular-value conditioning.  The displayed
nonnegative square root is not permission to hide a materially negative eigenvalue or
an unresolved rank classification.  Before a verification run, the verifier shall
derive and test arithmetic-error bounds for its Gram formation and eigensolution,
separately from the retained native-decimal uncertainty bound.  Clamping a negative
eigenvalue to zero is admissible only when its arithmetic-error interval includes
zero.  A materially negative interval or a singular-value interval that straddles the
frozen rank threshold gives numerical inconclusiveness and stops progression; it is
not rounded into agreement with the producer.  Conditional arithmetic-precision
escalation follows the preceding amendment; an inconclusive result never passes merely
because escalation is optional in already decisive cases.  No new numeric verification tolerance
is selected by this protocol revision.

The verifier shall report Hermiticity, orthonormality, and eigenpair residuals, including
\begin{equation}
 \rho_j(\mathbf{k})=
 \left\|G(\mathbf{k})v_j(\mathbf{k})-
 \lambda_j(\mathbf{k})v_j(\mathbf{k})\right\|_2.
\end{equation}
It shall derive the deficient atomic-trial direction from the minimum-eigenvalue
coordinate and independently aggregate its $s$, $p$, and $s^*$ weights.  It shall also
verify that a unitary basis change within the retained Wannier subspace does not change
the overlap singular spectrum.  This invariance control is not a fitted
momentum-dependent alignment $Q(\mathbf{k})$.

\section{Scientific interpretation and gate}

CSD1-V passes only if source identities, the complete 216-point rank field, deficient
point inventory, native-bound classification, deficient-direction weights, residuals,
and unitary-invariance controls agree within prospectively declared bounds.  A mismatch,
source-identity failure, numerical inconclusiveness, or independence-contract violation
produces an unresolved CSD1 disposition and stops the sequence.  Agreement verifies the
rank diagnosis under the declared finite-precision contract; it does not identify the
symmetry cause and does not validate a Slater--Koster model.

\chapter{CSD2: symmetry-sector diagnosis}

\section{Process}

CSD2 begins only after CSD1-V passes.  It shall inventory the little-group operations
available at each deficient point and authenticate their action on both the parent-space
trial columns and retained Wannier frame.  No symmetry representation may be inferred
from filenames, nominal orbital labels, energy degeneracy, or geometric expectation.  If
the retained artifacts do not determine the necessary actions and correlations, CSD2
shall record \texttt{INSUFFICIENT\_RETAINED\_SYMMETRY\_DATA} and stop.

\section{Mathematics}

For a deficient point $\mathbf{k}_*$ with little group
$\mathcal{G}_{\mathbf{k}_*}$, let $D_C(g)$ and $D_W(g)$ denote authenticated actions on
the trial and retained represented subspaces.  The overlap is symmetry-compatible only
when the corresponding intertwining relation is satisfied within a declared residual:
\begin{equation}
 D_W(g)^{\dagger}B(\mathbf{k}_*)D_C(g)=B(\mathbf{k}_*)
 \qquad (g\in\mathcal{G}_{\mathbf{k}_*}).
 \label{eq:csd2-intertwiner}
\end{equation}
When authenticated characters and group closure are available, sector projectors may be
constructed as
\begin{equation}
 \Pi_{\alpha}=\frac{d_{\alpha}}{|\mathcal{G}_{\mathbf{k}_*}|}
 \sum_{g\in\mathcal{G}_{\mathbf{k}_*}}
 \chi_{\alpha}(g)^*D(g).
\end{equation}
CSD2 shall compare trial and retained sector dimensions and localize the rank loss only
when those projectors, their idempotence, orthogonality, completeness, and representation
residuals are authenticated.

\section{Scientific interpretation and gate}

CSD2 distinguishes a symmetry-sector mismatch from an unexplained orbital-subspace
rank loss.  Energy degeneracy alone is not a symmetry label, and a symmetry-sector
finding does not establish that the source is physically inadequate.  CSD2 may report a
supported sector mismatch, a supported absence of such a mismatch under the admitted
data, or insufficient retained data.  It does not select a successor branch.

\chapter{CSD3: scientific disposition and successor selection}

\section{Process}

CSD3 is a human scientific decision over the authenticated CSD0 record, the independently
verified CSD1 result, and the CSD2 disposition.  The decision record shall name the
scientific question, admissible branch, fixed source, state space, basis, gauge, energy
reference, units, geometry, prohibited fallbacks, and the evidence required to stop or
continue.  It shall not reinterpret a failed rank-ten gate as a successful fit.

\section{Decision alternatives}

The admissible alternatives are:
\begin{enumerate}
 \item retain the blocked S110 conclusion and perform no successor study;
 \item select a spectral-only feasibility branch with no operator-equivalence claim;
 \item select a global-unitary rank-ten branch with a prospectively frozen
   momentum-independent $Q\in U(10)$;
 \item select a separately declared rank-nine projected-source branch; or
 \item select a new atomic-semantic Wannier source requiring separate scientific review
   and separately authorized QE/Wannier90 execution.
\end{enumerate}
CSD3 may select at most one active successor question for the next protocol tranche.
Selection does not execute that branch.  Before any selected branch is activated, the
HF0--HF4 hopping-first prerequisites below must be completed and accepted.

\section{Scientific interpretation and gate}

CSD3 owns interpretation and branch choice; software does not.  A global-unitary study
asks whether one fixed basis transformation can align the represented operators.  A
spectral-only study asks a weaker eigenvalue question.  A rank-nine study changes the
represented state space.  A new-source study changes the physical construction.  These
questions are not interchangeable, and failure or success in one shall not be promoted
to another.

\chapter{Hopping-before-Slater--Koster dependency}

This calculation adopts the dependency logic of the
\href{../../../periodic/multiband/analysis-order.md}{periodic multiband analysis-order
record}: identify the common state space, verify the complete Fourier representation,
measure withheld interpolation before support removal, characterize hopping locality
descriptively, make any irreversible support decision, report its separate
consequences, and only then attempt a constrained parameterization.  That program-level
record donates order and vocabulary only.  It transfers no silicon geometry, gauge,
threshold, shell decision, evidence status, or validation claim.

The historical chronology is retained rather than rewritten.  S30 characterized the
complete operator, S90 froze the complete shell family, S100 evaluated every shell,
S105 retained full support, and S110 then stopped before fitting because its proposed
atomic frame failed the common-space rank gate.  The reusable Level-1 Slater--Koster
software foundation was later implemented with synthetic tests only.  None of those
records is discarded or relabeled.  This protocol adds the prospective CSD0--CSD3
diagnosis and disposition sequence, followed by the HF0--HF4 hopping-first synthesis
gate, before any future spectral-only fit, global-unitary fit, rank-nine study, or
new-source Slater--Koster study.

\chapter{HF0: representation and authority gate}

\section{Process}

The hopping study shall authenticate the exact S30 full source, its direct lattice,
Wannier90 degeneracy convention, element-resolved WS shifts, gauge matrices, energy
reference, units, and the retained S90, S100, S105, and S110 dispositions.  Its source is
$H_{\mathrm{W}}$; under S105 the selected $H_{\mathrm{T}}$ equals
$H_{\mathrm{W}}$.  The revealed 408-point parent data may be used only as labeled
development-oracle evidence, and the 402-point parent values remain unobserved.  A
failed identity, basis, gauge, or convention check yields an unresolved hopping study
rather than a guessed conversion.

\section{Mathematics}

HF0 fixes the represented Fourier pair
\begin{equation}
 H_{\mathrm{W}}(\mathbf{k})=
 \sum_{\mathbf{R}}H_{\mathrm{W}}(\mathbf{R})
 \exp\!\left(+2\pi i\,\mathbf{k}\cdot\mathbf{R}\right),
 \label{eq:hf0-fourier-pair}
\end{equation}
with the exact ordinary-cell displacement, degeneracy convention, basis order, gauge,
and reciprocal-coordinate convention attached to every coefficient.  Fractional orbital
centers, when used by a later admitted model, enter through a separate diagonal gauge
transformation and do not redefine $\mathbf{R}$.

\section{Scientific interpretation and gate}

HF0 establishes the identity of the fixed source and the meaning of its coefficients.
It does not establish locality, atomic semantics, a finite-range model, or physical
accuracy.  HF0 passes only when the source and every convention are exact and
order-sensitive.

\chapter{HF1: complete hopping representation before support removal}

\section{Process and mathematics}

Analyze the complete set of 383 element-resolved WS representatives; do not preselect a
radius, retain only the largest entries, or infer a Slater--Koster shell map.  For every
representative and every complete physical-radius shell, retain the integer lattice
translation, element-resolved WS displacement convention, multiplicity or degeneracy,
pair partner, and matrix spectral, Frobenius, and maximum-entry norms.  Distinguish the
Bravais translation from any center-to-center physical displacement.  Verify
$R\leftrightarrow -R$ closure, Hermitian pairing, exact source-mesh Fourier
reconstruction within the native-serialization bound, and the applicable finite
Parseval identity under the declared convention.

These matrices live in the fixed rank-ten Wannier gauge.  Individual entries and an
entry-wise ``orbital hopping'' interpretation are gauge dependent.  The failed S110
TB-anchored alignment means that the existing source does not acquire two-site
$(s,p_x,p_y,p_z,s^*)$ semantics merely because its matrix rank is ten.  Site, orbital,
$\sigma$, or $\pi$ labels are forbidden unless a separately admitted common-space map
supports them.  Until then, report gauge-fixed matrix blocks, complete-shell norms, and
gauge-invariant spectra with the limitation explicit.

\section{Scientific interpretation and gate}

HF1 establishes a complete, convention-authenticated representation of the fixed
source.  Decay of shell norms is descriptive evidence in the retained Wannier gauge.
It neither removes support nor grants the matrix entries site-orbital meaning.  Missing
representatives, failed reverse closure, failed reconstruction, or failed Parseval
controls stop the sequence.

\chapter{HF2: development evaluation and descriptive locality}

\section{Process}

Before removing support, evaluate the full $H_{\mathrm{W}}(\mathbf{k})$ and every
complete-shell candidate on the exact 408-point development inventory.  Measure each
candidate relative to full $H_{\mathrm{W}}$ in the same rank-ten Wannier gauge.  This is
fixed-gauge finite enumeration: it does not modify, wrap, rerun, or reimplement Wannier90
optimization.  The C6 construction mesh remains reconstruction evidence, not the development evaluation
inventory.  The frozen 402-point geometry and parent values are forbidden inputs.  A
revealed 408-point direct-parent comparison may be displayed only as a separate
evaluation overlay and may not affect candidates, metrics, frontier membership,
truncation selection, fitting, or thresholds.

\section{Mathematics and channel separation}

Report cumulative retained norm, excluded-tail norm, the spectral triangle tail bound,
the Frobenius/Parseval tail, shellwise anisotropy, cancellation and nonmonotonicity, and
stability of spectra, valleys, curvature, masses, and named edge observables as complete
shells are traversed.  Tail norms and bounds remain diagnostics rather than substitutes
for measured channel errors.  Descriptive locality does not itself remove a hopping and
does not establish exponential decay, a continuum bound, an atomic-orbital
interpretation, or a finite-range Slater--Koster model.

For each exact scalar metric, classify a separate descriptive Pareto frontier between
retained support cost and that metric.  The required channels are operator error,
ordered-spectrum error, gap change, valley displacement and energy change, and each
named curvature/mass measure, all relative to full $H_{\mathrm{W}}$ on the 408-point
inventory.  Do not scalarize channels, average ranks, or select a compromise shell.
Undefined valley or mass observables retain typed inapplicable outcomes rather than
numeric sentinels.  Exact identities, formulas, controls, outputs, and typed dispositions
are fixed by
\path{calculations/Si/bulk/analysis/c6-b20-w21-hopping-pareto-synthesis/README.md}.
The selected numerical policy uses exact integer support cost and conservative
uncertainty intervals.  One candidate robustly dominates another only when it is no
worse under interval ordering and strictly better in at least one cost/error coordinate;
positive interval overlaps remain unresolved and incomparable.

\section{Scientific interpretation and gate}

HF2 characterizes support/error tradeoffs relative to the fixed full source.  A frontier
is not an acceptance rule, a hidden-parent predictor, or a scalar model-selection
objective.  HF2 passes when every required channel has an exact or typed unresolved
outcome and the 402-point confirmation boundary and reserved 108/108 fitting split have
not been consumed.

\chapter{HF3: support consequence and current disposition}

\section{Process and mathematics}

Only after HF1 and HF2 may support be removed.  The existing S90/S100/S105 branch is the
retained support study for this source: all 18 complete shells were evaluated, no
prospective nontrivial acceptance threshold existed, and full support was retained.
Thus the current controlled chain has
\begin{equation}
 H_{\mathrm{T}}=H_{\mathrm{W}},\qquad
 H_{\mathrm{W}}-H_{\mathrm{T}}=0.
\end{equation}
The other 17 shells remain descriptive evidence and are not retroactively rejected.
Any future nontrivial support choice requires a new prospective disposition and cannot
reuse parent numerical-convergence tolerances as model-reduction tolerances.

\section{Scientific interpretation and gate}

For the present source, HF3 preserves the full operator and therefore introduces no
spatial-truncation error.  This identity does not imply that the source is physically
accurate or that a short Slater--Koster model exists.  A future support change is a new
scientific decision, not a reinterpretation of S105.

\chapter{HF4: handoff to constrained parameterization}

\section{Process}

A hopping-first synthesis shall freeze its inputs, complete-shell inventory,
channel-specific descriptive frontiers, limitations, and unchanged S105 full-support
disposition before a future Slater--Koster fit is authorized.  It shall not claim hidden
error prediction.  Hopping locality and Slater--Koster parameterization remain distinct:
small long-range coefficients do not prove that a chosen angular parameter family spans
the retained operator, and a short Slater--Koster model does not retroactively define
the source's gauge or orbital semantics.  Model classes, neighbor maps, alignment
family, fit and validation roles, optimizer budgets, and failure outcomes require their
own prospective contract.

\section{Scientific interpretation and gate}

HF4 freezes the hopping evidence supplied to the selected successor branch.  It keeps
model definition separate from fit, residual construction, sample policy, model
selection, evidence, and acceptance.  Completion of HF4 permits preparation of the
CSD3-selected branch; it does not authorize execution or establish model adequacy.

The already implemented reusable model-definition foundation is preserved as software
evidence, but it does not close HF0--HF4 and does not authorize a scientific fit.  After
HF4 is accepted, only the CSD3-selected branch remains available for preparation under
separate human authorization.

\chapter{Post-S110 successor-branch boundary}

The spatial-truncation disposition retained complete support, so the selected
$H_{\mathrm{T}}$ equals $H_{\mathrm{W}}$ and the selected spatial-truncation error is
identically zero.  The subsequent S110 attempt tested the declared ten-orbital
TB-anchored map
\begin{equation}
 B(\mathbf{k})=
 \left[U_{\mathrm{dis}}(\mathbf{k})U(\mathbf{k})\right]^{\dagger}
 A(\mathbf{k})T^{\dagger},
\end{equation}
where $T$ converts the eight native $sp^3$ hybrids and two radial $s$ projections to
two site-resolved $(s,p_x,p_y,p_z,s^*)$ blocks.  A unitary common-space identification
requires all ten singular values of $B(\mathbf{k})$ to be resolved at every construction
point.  Under the conservative native-format singular-value uncertainty bound
$5.890206739032098\times10^{-7}$, the resolved rank is nine at Gamma and at the ordered
point $(0,1/3,0)$.  Their minimum singular values are respectively
$5.362000386100462\times10^{-13}$ and
$2.0186562353866025\times10^{-7}$.  No fit ran and no Slater--Koster parameter vector
exists.  Thus S50 remains blocked under the original TB-anchored operator contract.

Stopping the Slater--Koster branch and retaining this rank failure is the default
disposition under that contract.  Continuation is permitted only as a separately named, separately authorized study.
CSD3 selects the scientific question; HF0--HF4 must then close before the selected
branch is activated.  The following chapters define the available branches.  They are
alternatives, not an automatic ladder.

\chapter{Reusable Slater--Koster model-definition foundation}

\section{Process, mathematics, and scientific boundary}

Before any successor fit, reusable software shall implement the spinless ten-orbital
$sp^3s^*$ Hamiltonian, exact geometry and phase conventions, first- and second-neighbor
shell maps, real parameter domains, Hermiticity, symmetry checks, and deterministic
serialization.  Campaign declarations and evidence remain calculation-owned.  The
initial nearest-neighbor model has on-site parameters $E_s$, $E_p$, and $E_{s^*}$ and
hoppings $V_{ss\sigma}^{(1)}$, $V_{sp\sigma}^{(1)}$,
$V_{pp\sigma}^{(1)}$, $V_{pp\pi}^{(1)}$, and
$V_{s^*p\sigma}^{(1)}$; first-neighbor $ss^*$ and $s^*s^*$ channels remain fixed to
zero.  The declared second-neighbor extension adds
$V_{ss\sigma}^{(2)}$, $V_{sp\sigma}^{(2)}$,
$V_{pp\sigma}^{(2)}$, $V_{pp\pi}^{(2)}$,
$V_{s^*p\sigma}^{(2)}$, $V_{ss^*\sigma}^{(2)}$, and
$V_{s^*s^*\sigma}^{(2)}$.  Synthetic software tests establish implementation behavior
only; this level performs no scientific fit and establishes no model adequacy.

\chapter{Spectral-only feasibility branch}

\section{Process and mathematics}

A spectral-only study may fit the ordered eigenvalues or prospectively declared
symmetry-aware spectral invariants of $H_{\mathrm{SK}}$ to those of $H_{\mathrm{T}}$.
The already frozen, disjoint 108/108 construction-grid fit and validation roles are
reserved exclusively for that future, separately authorized study; they are not consumed
or reinterpreted by the hopping synthesis.  The study shall record every band-matching
convention, fit the declared model classes separately, and validate held-out spectra,
gap, valleys, curvatures, and masses.  It may use neither the revealed 408-point parent
values nor the unobserved 402-point values as fitting objectives or selection criteria.

\section{Scientific interpretation and gate}

Because this study supplies no justified orbital map, it shall not evaluate or claim an
operator difference $H_{\mathrm{T}}-H_{\mathrm{SK}}$.  Its outputs are spectral and
observable discrepancies only.  A successful result establishes at most descriptive
spectral feasibility for the declared construction source; it does not satisfy S110's
aligned-operator contract.

\chapter{Global-unitary rank-ten operator branch}

\section{Process and mathematics}

The preferred no-calculator continuation for the full rank-ten operator question is a
new study that fits one $\mathbf{k}$-independent unitary $Q\in U(10)$ jointly with the
Slater--Koster parameters $\theta$:
\begin{equation}
 \min_{Q\in U(10),\,\theta}
 \sum_{\mathbf{k}\in K_{\mathrm{fit}}}
 \left\|
 H_{\mathrm{T}}(\mathbf{k})-
 QH_{\mathrm{SK}}(\mathbf{k};\theta)Q^{\dagger}
 \right\|_{F}^{2}.
\end{equation}
The unitary must be common to every k point and must preserve the frozen geometry,
energy reference, orbital order, and spin convention.  Before fitting, the study shall
freeze its unitary parameterization, deterministic initializations, alternating or
manifold optimizer, stopping rules, model-class order, held-out validation, and failure
disposition.  With $Q$ fixed, linear parameter solves may be used; with $\theta$ fixed,
optimization of $Q$ remains a constrained nonlinear problem.  A spectral-only solution
may be declared as one initialization only if that role is recorded before the joint
fit.

\section{Scientific interpretation and gate}

The result shall report start-to-start stability, parameter identifiability, unitary
non-uniqueness, fit and validation operator residuals, spectra, and source-relative
observables.  A fit is not accepted merely because one initialization reaches a small
training residual.  This level changes the alignment method from TB-anchored
identification to learned global identification, but, if its gates pass, it retains a
well-defined rank-ten $H_{\mathrm{T}}-H_{\mathrm{SK}}$ operator channel.

\chapter{Rank-nine projected-source branch}

\section{Process and mathematics}

A rank-nine study changes the reduction chain.  It shall first define and freeze a
nine-dimensional source map $W_9(\mathbf{k})$ and then distinguish
\begin{equation}
 H_{\mathrm{T}}(\mathbf{k})
 \longrightarrow
 H_9(\mathbf{k})=W_9(\mathbf{k})^{\dagger}
 H_{\mathrm{T}}(\mathbf{k})W_9(\mathbf{k})
 \longrightarrow H_{\mathrm{SK}}^{(9)}(\mathbf{k}).
\end{equation}
The $H_{\mathrm{T}}\rightarrow H_9$ target-space error is separate from the subsequent
parameterization error and cannot be hidden inside it.  Eligibility requires resolved
rank across the full construction grid, smooth and periodic $W_9(\mathbf{k})$, symmetry
covariance, a fixed energy reference, Hermiticity, and acceptable real-space locality
after transformation.  The unresolved Gamma direction alone does not prove that one
globally smooth and symmetry-compatible rank-nine subspace exists.  This level cannot
be described as the original rank-ten $sp^3s^*$ reduction.

\section{Scientific interpretation and gate}

A successful rank-nine construction would establish only a separately defined
nine-dimensional source and parameterization question.  It cannot satisfy the failed
rank-ten gate and cannot conceal the $H_{\mathrm{T}}\rightarrow H_9$ target-space
change.

\chapter{New atomic-semantic Wannier-source branch}

\section{Process and scientific boundary}

If site-resolved atomic $sp^3s^*$ interpretation is essential, a new target-space and
projection construction may be designed to make the ten-orbital trial frame full rank.
It defines a new source candidate rather than repairing the retained $H_{\mathrm{W}}$.
Fresh preprocessing, conversion, and localization each require their own bounded
calculator authorization.  Any eligible new operator must repeat source
characterization, hopping-first Pareto characterization, spatial-reduction disposition,
alignment, fitting, candidate freeze, and confirmation under new identities.  The
existing source and its failures remain immutable evidence.  No $C6$ band-count
extension, $C8$, or $C12$ occurrence follows automatically from this option.

\chapter{Successor regularization and progression rules}

Singular-value flooring, pseudoinverse alignment, or a singular nonorthogonal metric may
be implemented only as a labeled sensitivity experiment.  It shall not be accepted as
the operator reduction because its result depends on a regularization coordinate at the
point where state-space identification failed.

The mandatory progression is CSD0, CSD1, CSD1-V, CSD2, CSD3, and HF0--HF4.  After
HF4, only the branch selected by CSD3 may be activated under separate authorization.
The spectral-only, global-unitary, rank-nine, and new-source branches are not automatic
successors to one another.  Failure or non-identifiability at the
hopping/common-space or global-unitary level stops its dependent route.  Rank-nine
projection and a new Wannier source are separate research questions, not retries.  Every
reduction fit and model-class decision in the selected branch shall use the declared construction
source rather than the revealed 408-point parent values.  Those values may be displayed
only as a separate development evaluation overlay and may not calibrate or select a
model.  Every level must preserve the confirmation-only 402-point boundary, record any
use of construction-source information in a later initialization, and freeze all
hopping results, reductions, diagnostics, fits, and evaluation definitions before
confirmation-parent values are observed.

\chapter{Construction-mesh successor recommendation}

The negative parent-fidelity result ends optimizer tuning intended to repair that
candidate retrospectively.  It shall not be repaired by tolerance changes, additional
localization iterations, or mixing optimizer changes with a scientific-construction
revision.  It does not block controlled component-operator construction or
Slater--Koster reduction from the already converged and structurally eligible full
operator.  Such reductions inherit the source operator's physical-fidelity limitation,
and their reduction diagnostics and reduction errors shall be reported separately from
that limitation.

The first successor shall isolate reciprocal construction-mesh density.  Candidate
$C_6$ uses a Gamma-origin $6\times6\times6$ ordered mesh while holding fixed the accepted
$S_1$ structure and parent, QE 42/336~Ry cutoffs, fixed occupations,
\texttt{nosym} and \texttt{noinv} set true, \texttt{conv\_thr} set to
$1.0\times10^{-12}$~Ry, full-accuracy Davidson diagonalization, 12 parent bands,
rank 10, non-spinorial character, outer window
$[-5.85,19.02]$~eV, frozen window $[-5.85,8.00]$~eV, and the two projection lines
\texttt{Si:sp3} and \texttt{Si:l=0:r=2}.  The localization controls also remain
\texttt{dis\_num\_iter} 2000, \texttt{dis\_conv\_tol} $1.0\times10^{-10}$,
\texttt{dis\_conv\_window} 3, \texttt{dis\_mix\_ratio} 0.2,
\texttt{num\_iter} 4000, \texttt{conv\_tol} $2.0\times10^{-9}$~\AA$^2$,
\texttt{conv\_window} 6, enabled \texttt{precond}, \texttt{num\_cg\_steps} 200,
native-default \texttt{trial\_step} 2.0, unset \texttt{fixed\_step}, and disabled
\texttt{guiding\_centres}.

The $C_6$ route consists of distinct NSCF, Wannier90 preprocessing,
\texttt{pw2wannier90.x}, and localization occurrences.  Each calculator invocation
requires its own exact preflight and bounded authorization.  No occurrence in this
protocol is automatically authorized.

The first exact $C_6$ localization occurrence retained
\texttt{dis\_num\_iter}=2000 and completed natively, but exhausted that budget at an
absolute fractional change of $4.679\times10^{-8}$ against the unchanged dimensionless
$1.0\times10^{-10}$ criterion.  Wannierisation subsequently converged at iteration 316,
but cannot repair the unconverged subspace.  This is an optimizer-budget failure and
produces no provider-eligible operator.  The monotonically decreasing changes over
iterations 1800--2000 support preparation of one fresh bounded eligibility test with
\texttt{dis\_num\_iter}=4000.  Every target-space coordinate and all other optimizer
controls remain fixed, and no failed checkpoint or output may be consumed.  This budget
extension is an explicit optimizer-coordinate amendment, not evidence that the original
occurrence converged or that either local optimization reaches a global minimum.  The
fresh occurrence subsequently converged Disentanglement at iteration 3068, but
Wannierisation did not meet its six-consecutive-step spread-change gate by iteration
4000.  Its final absolute spread change was $1.26\times10^{-5}$~\AA$^2$ against the
unchanged $2.0\times10^{-9}$~\AA$^2$ criterion, and the late trajectory remained
nonmonotonic with large excursions.  The provider therefore rejected the output.  The
prospective stop rule now applies: no further automatic budget extension follows and
$C_8$ shall not be selected merely because $C_6$ lacks an eligible operator.

\section{Target-space candidate \texorpdfstring{$C_{6\text{-}T1}$}{C6-T1}}

The authenticated $C_6$ parent eigenvalues span
$[-5.8341319732,20.8036050385]$~eV.  The inherited outer window
$[-5.85,19.02]$~eV contains 12 parent states at 142 construction points, 11 at 62,
and only the rank-minimum 10 at 12; 86 states across 74 points lie above its maximum.
The fixed numerical window therefore does not define a mesh-stable parent-band
inventory.  This observation does not prove the cause of either optimizer failure, but
it ends pure mesh-only attribution and motivates a separate target-space candidate.

Candidate $C_{6\text{-}T1}$ changes only the outer window to
$[-6.0,21.0]$~eV.  On the authenticated 216-point interface this defines the outer
subspace as the span of all 12 computed parent states at every construction point, with
finite lower and upper margins of 0.1658680268 and 0.1963949615~eV.  The frozen window,
rank, projections, mesh, source order, energy reference, interface bytes, and all
optimizer controls remain fixed, including \texttt{dis\_num\_iter}=4000 and
\texttt{num\_iter}=4000.  The occurrence shall start fresh and consume no checkpoint or
failed localization output.  Because the window does not change NNKP topology or the
QE-generated EIG, AMN, and MMN data, those authenticated interface artifacts may be
reused without new preprocessing or conversion.

$C_{6\text{-}T1}$ is not a continuation of the one-variable mesh test.  Its prepared
localization request remains unexecuted and is not selected as the next occurrence after
the parent-band review below.

\section{AiiDA-moderate parent-band candidate
\texorpdfstring{$C_{6\text{-}B20}$}{C6-B20}}

The next selected target-space stage adopts the band-count heuristic, but not the
complete protocol, from \texttt{aiida-wannier90-workflows} release 3.1.0 at pinned
commit \texttt{1c0115cb7c77}.  The pinned moderate
workflow sets \texttt{num\_bands\_factor}=2.0, and its estimator takes the maximum of
the factor-scaled occupied-state and projection counts and the corresponding counts
plus four.\footnote{\href{https://github.com/aiidateam/aiida-wannier90-workflows/blob/1c0115cb7c77541a0bd549bc05005df435d651b6/src/aiida_wannier90_workflows/workflows/protocols/base/wannier90.yaml\#L6-L25}{Pinned moderate-protocol source};
\href{https://github.com/aiidateam/aiida-wannier90-workflows/blob/1c0115cb7c77541a0bd549bc05005df435d651b6/src/aiida_wannier90_workflows/utils/pseudo/__init__.py\#L292-L342}{pinned band-estimator implementation}.}
For the non-spin-polarized two-atom cell, eight pseudopotential valence electrons, and
ten declared trial projections, this gives
\begin{equation}
  N_{\mathrm{bands}}=\max(8,8,20,14)=20.
\end{equation}
For this parent-band decision, the calculation-owned adaptation cites the AiiDA
workflow only for that factor-two inventory rule.  The separately declared precise
convergence adaptation controls prospective localization and does not alter the band
count.  Neither adaptation adopts the workflow's default PBEsol pseudopotential,
derives projections through its pseudo machinery, or claims that AiiDA executed or
validated this campaign.

Candidate $C_{6\text{-}B20}$ changes only the NSCF parent-band count from 12 to 20.
The accepted structure and parent, PBE pseudopotential, 42/336~Ry cutoffs, ordered
Gamma-origin $6\times6\times6$ mesh, occupations, $10^{-12}$~Ry electronic threshold,
full-accuracy Davidson solver, and symmetry controls remain fixed.  Rank 10 remains the
prospective represented dimension; 20 calculated bands do not imply rank 20.

One separately authorized $C_{6\text{-}B20}$ NSCF completed successfully.  Its QEXSD
contains 216 points and 20 finite bands spanning $[-5.8341319732,30.2720123939]$~eV,
and its saved state verifies.  The ordered points and weights match the earlier
12-band $C_6$ occurrence exactly, while its first 12 eigenvalues agree within
$4.14\times10^{-11}$~eV.  Under the inherited $[-5.85,19.02]$~eV window, the enlarged
inventory contains 10--14 states per point with histogram
$\{10\!:\!12,11\!:\!62,12\!:\!96,13\!:\!24,14\!:\!22\}$.  Under the prepared
$[-6.0,21.0]$~eV window it contains 12--14 states with histogram
$\{12\!:\!48,13\!:\!108,14\!:\!60\}$.  Bands 13 and 14 therefore enter both
investigated windows; band 15 remains above 21~eV everywhere.

These observations authenticate the broader parent inventory.  The subsequent
calculation-owned target-space decision selects candidate
$C_{6\text{-}B20\text{-}W21}$ with outer window $[-6.0,21.0]$~eV.  This places the
upper boundary 0.1964~eV above the observed band-12 maximum and 0.2572~eV below the
observed band-15 minimum, while admitting 12--14 states and therefore at least two
outer-space dimensions beyond rank 10 at every point.  The frozen window, rank, and
projection family remain unchanged.  This selection is an optimizer-eligibility
target-space redesign, not evidence of convergence, a global minimum, or scientific
adequacy.

Bands 13--20 require new interface artifacts; the existing 12-band NNKP, EIG, AMN, and
MMN data cannot be reused.  Fresh preprocessing, conversion, and localization remain
distinct occurrences with separate authorization.  The AiiDA precise factor of 3.0
would give 30 bands under the same adaptation, but $C_{6\text{-}B30}$ is a conditional
upper diagnostic and does not follow automatically from failure of
$C_{6\text{-}B20\text{-}W21}$.

The separately authorized preprocessing, conversion, and localization occurrences later
completed for $C_{6\text{-}B20\text{-}W21}$.  Native disentanglement and
Wannierisation converged, and no-calculator characterization established a finite,
Hermitian, source-reconstructing full untruncated operator.  It is therefore eligible as
the fixed source for controlled reduction.  Its revealed 408-point comparison passed
all six valley-position criteria but failed the gap criterion and 17 of 18
principal-mass criteria.  Those failures are development-oracle labels: they limit
physical-fidelity claims, but they neither revoke structural eligibility nor block the
controlled reduction and diagnostic study.

The revealed 408-point parent set is the development evaluation overlay, but it is no
longer withheld for $C_6$ or any later candidate.  The existing 5\% indirect-gap, 3\%
Gamma--X valley-position, and 8\% per-mass targets remain unchanged parent-evaluation
criteria.  They are not operational diagnostics and do not define candidates,
frontiers, truncation selection, fitting objectives, or thresholds.  The exact full
operator, reduction declarations, diagnostic rules, fitting conventions, and evaluation
definitions shall be frozen before any new parent values are observed.  Failure of a
development-oracle criterion
does not trigger an automatic $8\times8\times8$ successor.  A separately declared
$C_8$ mesh study is appropriate only for a separately stated source-fidelity question;
it is not a repair step required by the controlled-reduction protocol.  Rank, band
inventory, windows, projections, and mesh remain separate scientific coordinates and
shall not be changed together without an explicit new study.

A $12\times12\times12$ construction, denoted $C_{12}$, is a bounded upper mesh
diagnostic rather than an automatic successor.  The $C_4$, $C_6$, $C_8$, and $C_{12}$
meshes contain 64, 216, 512, and 1728 points, respectively, so $C_{12}$ is 27 times the
original source-point inventory.  It becomes scientifically eligible for a separate
preflight only when a completed $C_8$ review still supports the mesh-density hypothesis
and either the intended-use criteria remain unmet or material mesh sensitivity remains.
No direct jump from $C_4$ to $C_{12}$ is part of this protocol.

Passing the development and prospective confirmation criteria at $C_{12}$ would
establish adequacy only for that finite represented operator and declared intended use.
A stronger claim of construction-mesh convergence requires either comparison with a
finer declared mesh or a prospective real-space tail/aliasing bound showing that omitted
Fourier content cannot materially change the target observables.  Mesh size alone does
not supply that bound, and $C_{12}$ cannot repair an inadequate band inventory, rank,
window, or projection family.

A future confirmation set shall remain disjoint from the $C_4$, $C_6$, possible $C_8$,
and possible $C_{12}$ construction meshes and from the revealed 408-point set.  Its geometry is fixed
prospectively as follows.  Six symmetry-projected centers lie on the signed Cartesian
axes at the mean revealed parent-valley radius
$0.9629174643168844$~\AA$^{-1}$, equal to $0.8381689694043609$ of Gamma--X.  Around each
center, use the right-handed rotated Cartesian basis
\begin{equation}
 U=\begin{pmatrix}
 1/\sqrt{3} &  1/\sqrt{2} &  1/\sqrt{6} \\
 1/\sqrt{3} & -1/\sqrt{2} &  1/\sqrt{6} \\
 1/\sqrt{3} &  0           & -2/\sqrt{6}
 \end{pmatrix}.
\end{equation}
For each center, include the center, all six axial and twelve mixed-sign quadratic
points at each of $0.002$, $0.006$, and $0.012$~\AA$^{-1}$, and the twelve axial capture
points at $0.024$ and $0.048$~\AA$^{-1}$.  After periodic de-duplication the intended
inventory contains 402 points.  Exact fractional coordinates, construction-disjointness,
and point roles must be generated and authenticated before use.

This confirmation geometry is parent-informed because its centers use the already
revealed parent minima; it shall not be described as blind to geometry.  Its direct
parent eigenvalues nevertheless remain unobserved.  The geometry is confirmation-only:
it may not define or tune diagnostics, candidates, Pareto frontiers, truncation
selection, fitting, thresholds, or numerical comparison policies.  Reduction
definitions, operational diagnostics, fitting conventions, and evaluation criteria must
be frozen before a separately authorized direct-parent occurrence reveals those values.

\chapter{Controlled-reduction and confirmation sequence}

For the retained Gamma-origin $6\times6\times6$ source with 20 calculated parent
bands, retained rank ten, outer window $[-6.0,21.0]$~eV, frozen window
$[-5.85,8.0]$~eV, and the declared silicon $sp^3$ plus second radial-$s$ projection
family, post-S110 work shall first complete CSD0, CSD1, CSD1-V, CSD2, and CSD3 in that
order.  It shall then complete
HF0--HF4 before activating the single successor branch selected by CSD3.  A stopped,
unresolved, or insufficient-data disposition at any gate blocks every dependent stage;
there is no automatic fallback or branch substitution.

For $C_6$ and any separately approved source operator:

\begin{enumerate}
  \item verify native convergence, exact source-mesh reconstruction, finite dimensions,
    transformation identities, representative closure, and Hermiticity without using
    withheld parent values as an eligibility gate;
  \item freeze the full untruncated operator as $H_{\mathrm{W}}$ and declare its basis,
    gauge, energy reference, geometry, units, and known source-fidelity limitations;
  \item characterize the complete real-space hopping representation and its Fourier
    reconstruction without removing support or assigning unsupported atomic-orbital
    semantics;
  \item evaluate all complete shells relative to full $H_{\mathrm{W}}$ on the 408-point
    development inventory and retain separate channel-specific Pareto frontiers without
    scalarization or support selection;
  \item only then define and evaluate any spatial truncation, retaining shell and tail
    diagnostics and separate full-to-truncated consequences;
  \item only after the support disposition define a Slater--Koster parameterization,
    common-space map, model class, neighbor inventory, fit roles, and validation roles;
  \item retain separately hopping-locality evidence, support-removal error,
    parameterization error, spectral and observable stability, and every failed
    reduction;
  \item display revealed 408-point parent comparisons only as separate development
    evaluation overlays that cannot change frontier membership, support, fits, or
    thresholds;
  \item freeze the hopping results, reductions, diagnostic rules, fitting conventions,
    and evaluation criteria before observing confirmation parent values; and
  \item reveal the 402-point parent values only through a separately authorized
    occurrence, then report confirmation comparisons without refitting, recalibration,
    frontier revision, support reselection, or a hidden-error-prediction claim.
\end{enumerate}

Failure in any channel shall be attributed to the applicable parent-model,
discretization, subspace, gauge, interpolation, spatial-truncation, or parameterization
error class.  A parent-fidelity failure limits physical claims but does not erase the
controlled reduction.  Conversely, small internal reduction error does not repair or
hide a large parent-to-source discrepancy.  Passing software checks or native
convergence establishes operational eligibility only, not physical validation.

\chapter{Nonclaims}

This protocol does not certify a global minimum, physical accuracy, pseudopotential
transferability, mesh convergence, effective masses, locality, or an accepted
Slater--Koster model.  It does not claim that an operational diagnostic predicts,
calibrates, or equals parent error.  Pareto frontiers are descriptive internal
full-source tradeoff summaries; parent comparisons are separate evaluation overlays.
The protocol authorizes no calculator invocation, analysis execution, evidence
publication, release, or retrospective acceptance.

\end{document}
